%=======================================================================
%  CHAPTER 3 - ANALYSIS OF LTI SYSTEMS IN THE FREQUENCY DOMAIN  (6 hrs)
%  Sources: SPP Sir deck 4 "Transform Analysis of LTI" (40 slides,
%           exported from the legacy .ppt, Oppenheim ch5 material:
%           single pole/zero magnitude and phase, group delay,
%           multiple poles and zeros);
%           DSAP All Chapter Summary pp.50-61 (H(z) from an LCCDE, BIBO
%           via the ROC, two full pole-zero worked examples) and pp.85-87
%           (linear phase);
%           NAC Sir Ch-2 pp.17-24, pp.112-115 (generalized linear phase).
%           The Chapterwise set has NO file for this chapter.
%  Numericals live in ch3-num.tex.
%=======================================================================
\section{Analysis of LTI Systems in the Frequency Domain}

{\color{sub}\footnotesize 6 hrs \gap$\cdot$\gap \mk{56 questions across the 45 papers}
\gap$\cdot$\gap worth about \mk{8.4 of 80 marks} \gap$\cdot$\gap 2 syllabus sub-topics
\gap$\cdot$\gap \mk{6th of the 8 chapters} by weight: only Ch 2 and the EX 753 extras
are lighter.}

\lead{Read this first:}
\begin{itemize}
  \item \textbf{35 of the 56 questions are one question}, asked again and again with the
        numbers changed: \emph{plot the poles and zeros in the $z$-plane, then sketch the
        magnitude response, not to scale} \tS{35}. It is worth 6 to 10 marks and it is
        the single most valuable thing in this chapter.
  \item It arrives in two disguises. Either a \textbf{difference equation} is given and
        you must find $H(z)$ first, or the \textbf{poles and zeros are handed to you}
        directly, sometimes in polar form. The second half of the work is identical.
  \item The remaining 21 questions are short bookwork worth 2 to 4 marks each:
        causality and stability from the impulse response and the ROC \tF{7}, the LCCDE
        and its system function \tF{6}, FIR against IIR \tP{3}, linear phase \tP{3},
        zero-input and zero-state response \tP{2}.
  \item \mk{Sketching does not mean guessing.} Three numbers, $|H|$ at $\omega=0$,
        $\pi/2$ and $\pi$, plus the pole and zero angles, fix the shape completely.
        \S3.3 gives the recipe.
  \item The cheapest check in the chapter: $H(e^{j0})=H(1)$ is the sum of all
        coefficients and $H(e^{j\pi})=H(-1)$ is their alternating sum. Both are one
        line of arithmetic, and both are marks.
\end{itemize}

\hr

\T{3.1 \quad Frequency response of an LTI system}

\Q{Derive the expression for the frequency response of an LTI system. Show that a
complex exponential is an eigenfunction of an LTI system}{\m{6}}
{\color{sub}\footnotesize \tP{2} \yr{\textbf{69 Bh}, \textbf{67 Mng}}
\gap$\cdot$\gap the opening paragraph of most answers in this chapter}

\sbsr{0.54}{%
\lead{The eigenfunction property}
Feed a complex exponential into an LTI system with impulse response $h[n]$ and convolve:
\[ y[n]=\sum_{k=-\infty}^{\infty}h[k]\,e^{j\omega(n-k)}
       =e^{j\omega n}\underbrace{\sum_{k=-\infty}^{\infty}h[k]e^{-j\omega k}}_{H(e^{j\omega})} \]
\begin{itemize}
  \item The input comes out unchanged in \emph{shape}, scaled by a complex number. So
        $e^{j\omega n}$ is an \textbf{eigenfunction} of every LTI system and
        $H(e^{j\omega})$ is the \textbf{eigenvalue}.
  \item $H(e^{j\omega})$ is the \textbf{frequency response}: the DTFT of $h[n]$, and
        equally $H(z)$ evaluated on the unit circle.
  \item[] \[ H(e^{j\omega})=H(z)\Big|_{z=e^{j\omega}} \]
  \item \mk{This substitution is legal only if the ROC contains the unit circle.} If it
        does not, the system has no frequency response at all, however neatly $H(z)$
        factors. See \S3.4.
\end{itemize}}{%
\lead{Response to a real sinusoid}
For $x[n]=A\cos(\omega_0 n+\phi)$ and a real $h[n]$, superposing the two exponentials gives
\[ y[n]=A\,|H(e^{j\omega_0})|\cos\!\big(\omega_0 n+\phi+\angle H(e^{j\omega_0})\big) \]
\begin{itemize}
  \item The system can only do two things to a sinusoid: \textbf{scale it} by $|H|$ and
        \textbf{shift it} by $\angle H$. It can never change its frequency.
  \item That is the whole reason the magnitude and phase responses are worth plotting.
\end{itemize}
\lead{Three properties that halve the work}
\begin{itemize}
  \item $H(e^{j\omega})$ is \textbf{periodic in $\omega$ with period $2\pi$}, because
        $e^{j\omega n}$ is.
  \item For real $h[n]$, $|H|$ is \textbf{even} and $\angle H$ is \textbf{odd} in
        $\omega$.
  \item So \mk{everything is known from $0\le\omega\le\pi$}, and that is the only range
        any paper asks you to draw.
\end{itemize}}

\lead{Magnitude, phase and group delay}
\begin{itemize}
  \item Writing $H(e^{j\omega})=|H(e^{j\omega})|\,e^{j\arg H(e^{j\omega})}$, the three
        standard plots are the \textbf{magnitude response} $|H|$, the \textbf{phase
        response} $\arg H$, and the \textbf{group delay}
  \item[] \[ \operatorname{grd}H(e^{j\omega})=-\frac{d}{d\omega}\arg H(e^{j\omega}) \]
  \item Group delay is the delay the system applies to the envelope of a narrow band of
        frequencies. \mk{Constant group delay means linear phase means no phase
        distortion}: see \S3.6.
  \item Log magnitude is quoted in dB as $20\log_{10}|H(e^{j\omega})|$.
\end{itemize}

\hr

\T{3.2 \quad The LCCDE and its system function}

\Q{State the linear constant coefficient difference equation and the corresponding
system function. Why do we need a difference equation?}{\m{3}\m{2}}
{\color{sub}\footnotesize \tF{6} \yr{73 Shr \m{3}, 71 Shr \m{3+2},
\textbf{69 Ch} \m{2}, \textbf{\texttt{69 Bh}} \m{2}, \textbf{73 Ch} \m{2}}
\gap$\cdot$\gap pure bookwork, \mk{same answer every year}}

\sbsr{0.5}{%
\lead{The equation}
An $N$th-order linear constant-coefficient difference equation relates the output samples
to the input samples with coefficients that do not change with $n$:
\[ \sum_{k=0}^{N}a_k\,y[n-k]=\sum_{k=0}^{M}b_k\,x[n-k] \]
\begin{itemize}
  \item \textbf{Linear}: no products or powers of $x$ or $y$, no $\sin y[n]$.
  \item \textbf{Constant coefficient}: $a_k$ and $b_k$ are numbers, not functions of $n$.
        This is what makes the system time-invariant.
  \item \textbf{Order} $N$: how many past outputs are fed back.
  \item Normalising $a_0=1$ gives the recursive computing form, which is how it is
        actually run on hardware:
  \item[] \[ y[n]=\sum_{k=0}^{M}b_k x[n-k]-\sum_{k=1}^{N}a_k y[n-k] \]
\end{itemize}}{%
\lead{The system function}
Z-transform both sides, using linearity and the time-shift property
$y[n-k]\leftrightarrow z^{-k}Y(z)$:
\[ Y(z)\sum_{k=0}^{N}a_k z^{-k}=X(z)\sum_{k=0}^{M}b_k z^{-k} \]
\[ H(z)=\frac{Y(z)}{X(z)}
      =\frac{\displaystyle\sum_{k=0}^{M}b_k z^{-k}}{\displaystyle\sum_{k=0}^{N}a_k z^{-k}}
      =\frac{B(z)}{A(z)} \]
\begin{itemize}
  \item \mk{An LCCDE always gives a rational $H(z)$}, and every rational $H(z)$ comes from
        some LCCDE. The two descriptions are the same object.
  \item \textbf{Zeros} come from the $x$ side (feedforward), \textbf{poles} from the $y$
        side (feedback). Read them straight off the two coefficient lists.
  \item Multiply top and bottom by $z^{N}$ to clear the negative powers before factoring.
\end{itemize}}

\lead{Why we need a difference equation \mk{(a 2 mark question, asked five times)}}
\begin{itemize}
  \item It is a \textbf{complete description} of an LTI system in the time domain, finite
        and exact, whereas $h[n]$ may be an infinitely long list.
  \item It is \textbf{recursive}: each output uses a handful of stored past values, so an
        IIR response of infinite length costs finite memory and finite arithmetic.
  \item It maps \textbf{directly onto hardware}: every term is one delay, one multiplier
        and one adder, which is exactly what a realisation structure draws (Ch 4).
  \item It gives \textbf{$H(z)$ by inspection}, so the frequency response, the poles, the
        zeros and the stability all follow from the same few coefficients.
  \item It carries \textbf{initial conditions}, which is what separates the zero-input
        from the zero-state response.
\end{itemize}

{\color{sub}\footnotesize A degree count worth remembering: if $M<N$, then $B(z)$ written
in positive powers is short by $N-M$, so \mk{$H(z)$ has $N-M$ zeros at the origin}. They
are real zeros and they belong on the plot, but being at $z=0$ they are equidistant from
every point of the unit circle and so \mk{contribute nothing to the magnitude shape}.}

\hr

\T{3.3 \quad Pole-zero plot to magnitude response}

\Q{Plot the poles and zeros in the $z$-plane and sketch the magnitude response (not to
scale) of the system described by the given difference equation}{\m{3+7}\m{2+8}\m{6}}
{\color{sub}\footnotesize \tS{35} \yr{\textbf{80 Bh}, 81 Ba, 76 Ash,
\textbf{79 Bh}, 79 Ba, \textbf{78 Bh}, \textbf{81 Bh}, \textbf{\texttt{80 Ch}},
\textbf{\texttt{72 Ash}}, 74 Ash, 72 Ka, \textbf{72 Ch}, 82 Ba, \textbf{70 Ch},
\textbf{\texttt{71 Bh}}, \texttt{73 Ma}, \textbf{69 Ch}, \textbf{\texttt{76 Bh}},
\textbf{75 Ch}, \textbf{69 Bh}, \textbf{67 Mng}, \textbf{68 Bh}, 70 Asa, \texttt{70 Ma},
80 Ba, \textbf{82 Bh}, \textbf{76 Ch}, 75 Ash, \textbf{\texttt{81 Ch}},
\textbf{\texttt{73 Bh}}, \textbf{74 Ch}, \textbf{\texttt{75 Bh}}, \textbf{73 Ch},
\textbf{71 Ch}}
\gap$\cdot$\gap \mk{worked in full in the numerical companion, \S1 and \S2}}

\sbsr{0.46}{%
\lead{The geometric idea, which is the whole method}
Factor $H(z)$ into its zeros $z_k$ and poles $p_k$:
\[ H(z)=K\,\frac{\prod_{k=1}^{M}(z-z_k)}{\prod_{k=1}^{N}(z-p_k)} \]
Now put $z=e^{j\omega}$, a point walking round the unit circle. Taking magnitudes turns
the product into distances:
\[ |H(e^{j\omega})|=|K|\,
   \frac{\prod_k\,\big|e^{j\omega}-z_k\big|}{\prod_k\,\big|e^{j\omega}-p_k\big|} \]
\begin{itemize}
  \item Every factor is the \textbf{length of the vector} drawn from a zero or a pole to
        the point $e^{j\omega}$.
  \item \mk{Zero distances multiply on top, pole distances divide underneath.}
  \item So as $\omega$ sweeps $0$ to $\pi$: passing \textbf{near a pole} makes a
        denominator distance small, and $|H|$ \textbf{peaks}; passing \textbf{near a
        zero} makes a numerator distance small, and $|H|$ \textbf{dips}.
  \item A zero exactly \emph{on} the unit circle gives a distance of zero, so $|H|$ is
        \textbf{exactly zero} there: a perfect null.
  \item The closer a pole sits to the circle, the taller and narrower its peak.
\end{itemize}}{%
\lead{The exam recipe}
\begin{enumerate}
  \item From the LCCDE write $H(z)=B(z)/A(z)$ in $z^{-1}$, then multiply top and bottom
        by $z^{N}$ so both are polynomials in $z$.
  \item \textbf{Factor}. Roots of the numerator are zeros, roots of the denominator are
        poles. Use the quadratic formula: these are almost always second order and the
        roots almost always come out complex conjugate.
  \item \textbf{Plot} the unit circle, mark each zero $\circ$ and each pole $\times$, and
        do not forget the zeros at the origin.
  \item \textbf{Evaluate three points}, which is all a not-to-scale sketch needs:
  \item[] \[ |H(e^{j0})|=\left|\frac{B(1)}{A(1)}\right|,\quad
             \left|H(e^{j\pi/2})\right|=\left|\frac{B(j)}{A(j)}\right|,\quad
             |H(e^{j\pi})|=\left|\frac{B(-1)}{A(-1)}\right| \]
  \item \textbf{Locate the peak} at the angle of the pole, and the dip at the angle of
        the zero. Join the points smoothly.
  \item State the character in words: low-pass, high-pass, band-pass or notch.
\end{enumerate}
{\color{sub}\footnotesize $B(1)/A(1)$ is just the sum of the coefficients, and
$B(-1)/A(-1)$ the alternating sum. \mk{Two lines of arithmetic for two anchor points of
the sketch.}}}

\lead{The algebraic form, for when a paper wants the expression and not only the picture}
Each first-order factor, written with the pole or zero at radius $r$ and angle $\theta$,
contributes
\[ \big|1-re^{j\theta}e^{-j\omega}\big|^{2}=1+r^{2}-2r\cos(\omega-\theta) \]
\begin{itemize}
  \item In dB: $\;20\log_{10}\big|1-re^{j\theta}e^{-j\omega}\big|
        =10\log_{10}\!\big(1+r^{2}-2r\cos(\omega-\theta)\big)$,
        \textbf{added} for a zero factor and \textbf{subtracted} for a pole factor.
  \item Phase: $\;\arg\big(1-re^{j\theta}e^{-j\omega}\big)
        =\arctan\dfrac{r\sin(\omega-\theta)}{1-r\cos(\omega-\theta)}$,
        again $+$ for a zero and $-$ for a pole.
  \item Group delay of the same factor:
        $\;\dfrac{r^{2}-r\cos(\omega-\theta)}{1+r^{2}-2r\cos(\omega-\theta)}$.
  \item \mk{The minimum of $1+r^2-2r\cos(\omega-\theta)$ is at $\omega=\theta$}, where it
        equals $(1-r)^2$. That is why a pole at angle $\theta$ puts the peak of $|H|$ at
        $\omega=\theta$, and why $r\to1$ makes it sharp.
\end{itemize}

\lead{Poles and zeros given in polar form}
{\color{sub}\footnotesize Several papers, and the local worked example the family is
copied from, specify a radius and an angle instead of a complex number. Convert first:}
\[ z=re^{j\theta}=r\cos\theta+jr\sin\theta \]
For example radius $0.9$ at angle $1.0376$ rad gives $0.45+j0.77$, and radius $0.892$ at
$2.5158$ rad gives $-0.728+j0.522$. \mk{Those are the exact numbers the papers reuse}, so
recognising them saves a conversion.

\hr

\T{3.4 \quad Causality and stability, from $h[n]$ and from the ROC}

\Q{Describe the stability and causality characteristics of an LTI system in terms of its
impulse response and the ROC of its transfer function, with suitable examples}{\m{4+3}\m{4}\m{2}\m{3}}
{\color{sub}\footnotesize \tF{7} \yr{76 Ash \m{4+3}, \textbf{72 Ch} \m{4},
\textbf{\texttt{81 Ch}} \m{2}, \textbf{\texttt{80 Ch}} \m{2}, \texttt{73 Ma} \m{2},
\textbf{\texttt{73 Bh}} \m{3}, \textbf{\texttt{77 Ch}} \m{3}}
\gap$\cdot$\gap \mk{learn the four-line table and the one-line rule}}

\sbs{%
\lead{In terms of the impulse response}
\par
\begin{tabular}{@{}L{2.1cm}L{5.3cm}@{}}
\toprule
\hrow \thead{Property} & \thead{Condition on $h[n]$} \\
Causal & $h[n]=0$ for all $n<0$ \\
Stable (BIBO) & $\displaystyle\sum_{n=-\infty}^{\infty}|h[n]|<\infty$ \\
\bottomrule
\end{tabular}
\par\smallskip
\begin{itemize}
  \item Causality says the output cannot anticipate the input: it is a statement about
        \emph{when} $h[n]$ is allowed to be non-zero.
  \item Stability says a bounded input always gives a bounded output, and the necessary
        and sufficient test is that $h[n]$ be \textbf{absolutely summable}.
  \item The two are independent. A system can be either, both or neither.
\end{itemize}}{%
\lead{In terms of the ROC}
\par
\begin{tabular}{@{}L{2.1cm}L{5.3cm}@{}}
\toprule
\hrow \thead{Property} & \thead{Condition on the ROC} \\
Causal & exterior of the outermost pole, and includes $z=\infty$ \\
Stable & the ROC \textbf{contains the unit circle} \\
Both & \textbf{all poles strictly inside} $|z|=1$ \\
\bottomrule
\end{tabular}
\par\smallskip
\begin{itemize}
  \item The last line is the one to quote: \mk{a causal system is stable if and only if
        every pole has $|p|<1$}.
  \item A pole exactly \emph{on} the circle is marginally stable, which is \textbf{not}
        BIBO stable: the frequency response blows up there.
  \item Stability is also why the frequency response may not exist. If the ROC misses the
        unit circle, $H(e^{j\omega})$ is undefined, so \mk{no magnitude response can be
        drawn at all}.
\end{itemize}}

\lead{One $H(z)$, three systems \mk{(the standard worked example, and a favourite trap)}}
Take $H(z)$ with poles at $z=\tfrac12$ and $z=3$. The algebra is identical in all three
cases; only the ROC differs, and the ROC is a separate piece of information the question
must supply.
\par\smallskip
\begin{tabular}{@{}L{2.2cm}L{4.4cm}L{2.0cm}L{2.4cm}@{}}
\toprule
\hrow \thead{ROC} & \thead{$h[n]$} & \thead{Causal?} & \thead{Stable?} \\
$\tfrac12<|z|<3$ & $(\tfrac12)^nu[n]-2(3)^nu[-n-1]$ & no & \textbf{yes} \\
$|z|>3$ & $(\tfrac12)^nu[n]+2(3)^nu[n]$ & \textbf{yes} & no \\
$|z|<\tfrac12$ & $(\tfrac12)^nu[-n-1]+2(3)^nu[-n-1]$ & no & no \\
\bottomrule
\end{tabular}
\par\smallskip
{\color{sub}\footnotesize The middle row is the point: \mk{a causal system with a pole
outside the unit circle is unstable}, no matter how ordinary $H(z)$ looks. The top row is
the only stable reading, and it is not causal.}

\hr

\T{3.5 \quad FIR against IIR}

\Q{Differentiate between an FIR system and an IIR system}{\m{4}}
{\color{sub}\footnotesize \tP{3} \yr{80 Ba \m{4}, \textbf{\texttt{76 Bh}} \m{4},
\textbf{67 Mng}} \gap$\cdot$\gap a table answer, four or five rows is enough}

\par
\begin{tabular}{@{}L{3.4cm}L{5.2cm}L{5.2cm}@{}}
\toprule
\hrow \thead{} & \thead{FIR} & \thead{IIR} \\
Impulse response & finite length, $h[n]=0$ outside $0\le n\le M-1$ & infinite length, decays but never ends \\
Difference equation & non-recursive, output from inputs only & recursive, past outputs fed back \\
$H(z)$ & all zeros, poles only at the origin & has poles away from the origin \\
Stability & \textbf{always stable}, a finite sum is always bounded & stable only if every pole is inside $|z|=1$ \\
Linear phase & \textbf{easy}, exact, by making $h[n]$ symmetric & not achievable exactly \\
Order needed & high, for a sharp response & low, for the same sharpness \\
Memory, arithmetic & more & less \\
\bottomrule
\end{tabular}
\par\smallskip
\begin{itemize}
  \item The two rows examiners look for are \mk{FIR is always stable} and \mk{FIR can have
        exactly linear phase}. Those are the reasons to pay for its higher order.
  \item An FIR system's $H(z)=b_0+b_1z^{-1}+\dots+b_{M-1}z^{-(M-1)}$ written over $z^{M-1}$
        has all its poles at $z=0$, which is inside the unit circle: the stability claim
        needs no further argument.
\end{itemize}

\hr

\T{3.6 \quad Linear phase}

\Q{What is the condition satisfied by a linear phase FIR filter? Derive the frequency
response of a symmetric linear phase filter of length $M$, $M$ odd}{\m{2+4}\m{6}}
{\color{sub}\footnotesize \tP{3} \yr{70 Asa \m{2+4}, \textbf{67 Mng} \m{6},
\textbf{69 Bh} \m{6}} \gap$\cdot$\gap \mk{the only chapter source is FIR-design
material}: this topic is taught again in Ch 5}

\sbsr{0.5}{%
\lead{What linear phase means, and why it is wanted}
\begin{itemize}
  \item Phase is \textbf{linear} when $\arg H(e^{j\omega})=-\alpha\omega$, a straight line
        through the origin.
  \item Then the group delay $-\dfrac{d}{d\omega}\arg H=\alpha$ is \textbf{constant}:
        every frequency is held up by the same $\alpha$ samples.
  \item A common delay shifts the signal but does not deform it. \mk{Linear phase means
        no phase distortion}, which matters wherever shape carries the information:
        speech, ECG, data pulses, images.
  \item If the phase were curved, different frequency components would arrive with
        different delays and a pulse would spread out, even with a perfectly flat $|H|$.
  \item \textbf{Generalized} linear phase allows a constant offset,
        $\arg H=\beta-\alpha\omega$. The group delay is still constant, so it still does
        not distort; $\beta=\pm\pi/2$ is what an antisymmetric $h[n]$ produces.
\end{itemize}}{%
\lead{The condition}
An FIR filter of length $M$ has linear phase if its impulse response is symmetric or
antisymmetric about its midpoint:
\[ h[n]=\pm\,h[M-1-n],\qquad 0\le n\le M-1 \]
\begin{itemize}
  \item $+$ gives \textbf{symmetric} $h[n]$ and phase $-\alpha\omega$ with
        $\alpha=\dfrac{M-1}{2}$.
  \item $-$ gives \textbf{antisymmetric} $h[n]$ and phase $\dfrac{\pi}{2}-\alpha\omega$,
        with the same $\alpha$.
  \item \mk{$\alpha=(M-1)/2$ is the centre of the sequence}, and it is the delay in
        samples. For odd $M$ it is an integer; for even $M$ it is a half-integer, so the
        filter delays by half a sample.
  \item In the $z$-plane the condition forces zeros to occur in conjugate reciprocal
        groups: if $z_0$ is a zero, so are $z_0^{*}$, $1/z_0$ and $1/z_0^{*}$.
\end{itemize}}

\lead{Derivation for symmetric $h[n]$ with $M$ odd \mk{(the 6 mark answer)}}
\begin{itemize}
  \item Start from the definition and pull out the centre term at $n=\alpha=(M-1)/2$:
  \item[] \[ H(e^{j\omega})=\sum_{n=0}^{M-1}h[n]e^{-j\omega n} \]
  \item Substitute $n\to M-1-n$ in the sum and use $h[n]=h[M-1-n]$ to pair each term with
        its mirror image. Each pair contributes
  \item[] \[ h[n]\big(e^{-j\omega n}+e^{-j\omega(M-1-n)}\big)
             =2h[n]\,e^{-j\omega\alpha}\cos\!\big(\omega(\alpha-n)\big) \]
  \item Summing the $\alpha$ pairs and the lone centre term factors the whole response:
  \item[] \[ H(e^{j\omega})=e^{-j\omega\alpha}
             \underbrace{\left[h[\alpha]+2\sum_{n=0}^{\alpha-1}h[n]
             \cos\!\big(\omega(\alpha-n)\big)\right]}_{\text{real, call it }H_r(\omega)} \]
  \item The bracket is \textbf{purely real}, so the entire phase is the exponential in
        front:
  \item[] \[ |H(e^{j\omega})|=|H_r(\omega)|,\qquad
             \arg H(e^{j\omega})=-\alpha\omega\ \ (\text{plus }\pi\text{ wherever }H_r<0) \]
  \item \mk{That is the proof}: the phase is a straight line of slope $-\alpha$, so the
        group delay is $\alpha$ samples at every frequency.
\end{itemize}

{\color{sub}\footnotesize The antisymmetric case runs identically but the mirror pair
subtracts, giving $2jh[n]\sin(\omega(\alpha-n))$ and hence
$H=je^{-j\omega\alpha}H_i(\omega)$. The extra $j$ is the constant $\pi/2$ in the phase,
which is why an antisymmetric filter has \mk{generalized} linear phase rather than the
plain kind. See \S5 of the numerical companion for $h[n]=\{-1,0,1\}$, which is this
case.}

\hr

\T{3.7 \quad Last-minute recall}

\sbs{%
\lead{Say these without thinking}
\begin{itemize}
  \item $H(e^{j\omega})=H(z)|_{z=e^{j\omega}}$, \textbf{provided the ROC contains the
        unit circle}.
  \item $|H|=\dfrac{\text{product of distances from the zeros}}{\text{product of
        distances from the poles}}$.
  \item Pole near the circle $\rightarrow$ peak; zero near the circle $\rightarrow$ dip;
        zero on the circle $\rightarrow$ exact null.
  \item Peak sits at the \textbf{angle of the pole}, dip at the \textbf{angle of the
        zero}.
  \item $|H(e^{j0})|=|B(1)/A(1)|$, $\ |H(e^{j\pi})|=|B(-1)/A(-1)|$.
  \item Causal and stable $\Leftrightarrow$ \textbf{all poles inside} $|z|=1$.
  \item FIR: always stable, exactly linear phase possible, all poles at the origin.
  \item Linear phase $\Leftrightarrow$ $h[n]=\pm h[M-1-n]$, delay $\alpha=(M-1)/2$.
\end{itemize}}{%
\lead{Mistakes that cost marks}
\begin{itemize}
  \item \mk{Forgetting the zeros at the origin} when $M<N$. They belong on the plot even
        though they do not shape $|H|$.
  \item Drawing a magnitude response for a system whose ROC excludes the unit circle. If
        a pole lies outside and the system is stated causal, \mk{the correct answer is
        that no frequency response exists}.
  \item Reading the poles off the $x$ coefficients. \textbf{Poles come from the $y$
        side.}
  \item Losing the conjugate. A real difference equation can only have complex roots in
        \textbf{conjugate pairs}; a lone complex pole means an arithmetic slip.
  \item Sketching a peak at the pole's \emph{radius} instead of its \emph{angle}.
  \item Quoting $\omega$ in Hz. It is in radians per sample, and the plot runs $0$ to
        $\pi$.
\end{itemize}}
